\chapter{Stand der Forschung} \label{kap:3}

Das folgende Kapitel gibt einen thematisch strukturierten Überblick über den Forschungsstand zur \ac{KI}-gestützten Barrierefreiheitsanalyse. Darauf aufbauend werden die verbleibenden Forschungslücken herausgearbeitet, die den Ausgangspunkt der vorliegenden Arbeit bilden. Die relevanten Arbeiten lassen sich drei Entwicklungslinien zuordnen: (1) \ac{LLM}-basierte Erkennung von Accessibility-Verstößen, (2) \ac{LLM}-basierte Korrektur und Code-Generierung sowie (3) hybride Pipelines, die regelbasierte Tools mit \ac{LLM}-Analyse kombinieren.

\section{Vorgehen der Literaturrecherche}

Zur Einordnung des Forschungsstands wurde eine strukturierte Literaturrecherche in wissenschaftlichen Datenbanken durchgeführt (ACM Digital Library, IEEE Xplore, ScienceDirect, arXiv, ResearchGate).\footcite[Vgl.][]{WebsterWatson2002} Die Recherche erfolgte themenorientiert auf Basis von Suchbegriffen wie \glqq web accessibility\grqq, \glqq LLM accessibility\grqq, \glqq accessibility violation detection\grqq{} und \glqq axe-core\grqq; berücksichtigt wurden primär aktuelle Arbeiten mit Bezug zur automatisierten Erkennung oder Behebung von Barrierefreiheitsverstößen in Webanwendungen mit \ac{KI}-Assistenzsystemen. Die identifizierten Arbeiten wurden thematisch gruppiert, priorisiert und vergleichend ausgewertet.

\section{Von frühen \acs{KI}-Ansätzen zur \acs{LLM}-basierten Erkennung}

Delnevo et al. (2024) zeigen, dass GPT-3.5 einfache syntaktische Verstöße mit moderater Genauigkeit erkennt, bei semantischen Problemen jedoch inkonsistente Bewertungen liefert.\footcite[Vgl.][]{delnevo_interaction_2024} Bassi et al. (2025) vergleichen GPT-4o, GPT-3.5 und LLaMA~3.1 in einer zweistufigen Evaluation; GPT-4o identifiziert 74\% der Verstöße korrekt, während 15\% als Halluzinationen eingestuft werden.\footcite[Vgl.][S. 302]{bassi_assessment_2025} Eine nachgelagerte Chain-of-Verification-Stufe erhöht die Zuverlässigkeit deutlich; \ac{RAG} wird nur als zukünftige Erweiterung diskutiert.\footcite[Vgl.][S. 301]{bassi_assessment_2025}

\section{\acs{LLM}-basierte Korrektur und Prompt-Engineering}

Abu Doush und Kassem (2024) zeigen, dass strukturierte Prompts die Qualität \ac{LLM}-generierter barrierefreier \acs{HTML}-Alternativen erheblich verbessern, bei komplexen Interaktionsmustern aber weiterhin Schwächen bestehen.\footcite[Vgl.][S. 343]{doush_evaluating_2025} Guriță und Vatavu (2025) belegen, dass accessibility-orientierte Prompts paradoxerweise mehr Violations erzeugen können als neutrale -- erst iteratives Prompting über fünf Runden senkte den Violation Score von 0,84 auf 0,32.\footcite[Vgl.][S. 1000]{gurita_insights_2025}

Den aktuellen Stand der Technik repräsentiert \textfett{AccessGuru} von Fathallah et al. (2025).\footcite[Vgl.][]{fathallah_accessguru_2025} Das System kombiniert eine zweistufige Pipeline -- zunächst regelbasierte Detection via Playwright und \textfett{axe-core}, anschließend eine \ac{LLM}-basierte Korrektur durch GPT-4o mit Role-Play, Contextual und Metacognitive Prompting. Corrective Re-Prompting verbessert den Violation Score Decrease von 0,72 auf 0,84 bei einer semantischen Ähnlichkeit von 77\% zu menschlichen Referenzlösungen. Als offenes Problem benennen die Autoren die fehlende Verknüpfung korrigierter \acs{HTML}-Snippets mit dem Quellcode der geprüften Anwendung.

Fernández-Navarro und Chicano (2026) erweitern diesen Ansatz um eine direkte Quellcode-Korrektur und adressieren damit explizit die Lücke zwischen \ac{DOM}-Korrektur und entwicklerseitiger Umsetzung.\footcite[Vgl.][]{fernándeznavarro2026automatedllmbasedaccessibilityremediation} Ihr Tool injiziert Korrekturen auf Basis einer Quellcodeanalyse gezielt in tatsächliche Komponentendateien und erreicht in der Evaluation über öffentliche Websites und Open-Source-Angular-Projekte Remediation Rates von über 80\% bei vollständiger Build-Integrität. Der Ansatz beschränkt sich jedoch auf Angular und setzt auf cloudbasiertes GPT-4o; eine Übertragung auf React wird von den Autoren als kurzfristiger Folgeschritt benannt.\footcite[Vgl.][]{fernándeznavarro2026automatedllmbasedaccessibilityremediation}

\section{Hybride Pipelines und Ensemble-Testing}

Huang et al. entwickeln 2024 mit \textfett{ACCESS} ein System, das Prompt-Engineering gezielt zur automatisierten Barrierefreiheitskorrektur einsetzt und dabei verschiedene Strategien systematisch vergleicht.\footcite[Vgl.][S. 7,9]{huang_access_2024} Der Vergleich von ReAct, Chain-of-Thought und Few-Shot Prompting zeigt, dass ReAct mit einem Severity Score Reduction von 51,3\% alle anderen Strategien übertrifft -- ein Ergebnis, das die Bedeutung iterativer, werkzeuggestützter Reasoning-Strategien für diese Aufgabe unterstreicht.

He et al. präsentieren 2025 mit \textfett{GenA11y} das methodisch ausgefeilteste Detection-System des aktuellen Forschungsstands.\footcite[Vgl.][FSE101:20--21]{he_enhancing_2025} \textfett{GenA11y} extrahiert nach vollständigem clientseitigem Rendering den \ac{DOM} via Playwright und prüft ihn anschließend kriteriumsspezifisch mit GPT-4o. Eine Ablation Study belegt, dass weder \ac{DOM}-Extraktion noch optimiertes Prompting allein die Gesamtleistung von Precision 94,5\% und Recall 87,61\% erreichen.\footcite[Vgl.][FSE101:20]{he_enhancing_2025} \textfett{GenA11y} erkennt acht Verstoßtypen, die von keinem der verglichenen regelbasierten Tools gefunden werden.\footcite[Vgl.][FSE101:20]{he_enhancing_2025} Paternò et al. ergänzen dieses Bild mit \textfett{TeVal}, das semantische \ac{WCAG}-Kriterien durch Role-Play Prompting mit Few-Shot-Beispielen und Confidence Score bewertet und dabei eine Effektivität von 92,4\% erzielt.\footcite[Vgl.][]{paterno_how_2025}

Pool (2023) adressiert die Detection-Schicht mit \textfett{Testaro}, einem Ensemble aus acht Accessibility-Werkzeugen (u.\,a. axe-core, WAVE, QualWeb); die Analyse zeigt, dass jedes Tool exklusive Violations erkennt und die Werkzeuge damit komplementär wirken.\footcite[Vgl.][]{pool_testaro_2023}

\section{Forschungslücken und Einordnung der vorliegenden Arbeit} \label{Forschungslücken}

Die Analyse der bestehenden Forschungsarbeiten zeigt drei persistente Forschungslücken, die den Ausgangspunkt der vorliegenden Arbeit definieren.

Erstens berücksichtigt keiner der untersuchten Ansätze systematisch zustandsabhängige Barrieren, die erst durch Nutzerinteraktionen in React-basierten Single-Page Applications entstehen.\footcite[Vgl.][S. 110]{tafreshipour_ma11y_2024} Zwar setzen einzelne Systeme wie \textfett{AccessGuru} Playwright zur \ac{DOM}-Extraktion ein, jedoch beschränkt sich dies auf das initiale Rendering -- interaktionsbedingte Zustandswechsel werden nicht geprüft.\footcite[Vgl.][]{fathallah_accessguru_2025}

Zweitens verknüpft kein bekannter Ansatz Tool-Reports mit dem Quellcode der geprüften Anwendung. \textfett{AccessGuru} selbst benennt dies als offenes Problem: die generierten Korrekturen beziehen sich auf isolierte \ac{DOM}-Snippets und nicht auf die \ac{JSX}-Komponenten, in denen Entwicklerinnen und Entwickler tatsächlich eingreifen müssten.\footcite[Vgl.][]{fathallah_accessguru_2025} Die Lücke zwischen technischer Fehlererkennung und entwicklernaher Umsetzung bleibt damit bestehen.

Drittens adressiert keine der untersuchten Arbeiten die datenschutzrechtlichen Anforderungen öffentlicher Verwaltungen: alle Systeme setzen auf cloudbasierte \ac{LLM}-\ac{API}s (primär GPT-4o), was für die \ac{BITBW} und vergleichbare Institutionen ohne vertiefte datenschutzrechtliche Prüfung nicht nutzbar ist.\footcite[Vgl.][Art. 28, 44 ff.]{dsgvo2016}

Die Konzeptmatrix in Tabelle \ref{tab:konzeptmatrix_haupt} fasst die zentralen Merkmale der untersuchten Studien im Vergleich zum eigenen \ac{PoC} zusammen und macht die identifizierten Forschungslücken strukturell sichtbar.

\newcolumntype{A}{>{\raggedright\arraybackslash}p{4.3cm}} % Autor/Jahr
\newcolumntype{S}{>{\raggedright\arraybackslash}p{2.8cm}} % System
\newcolumntype{C}{>{\centering\arraybackslash}p{0.55cm}}  % Konzept-Spalten (alle gleich breit)

\newcommand{\lhdr}[2]{\parbox[c][32mm][c]{#1}{\centering\textbf{#2}}}
\newcommand{\hdr}[1]{\rotatebox{90}{\scriptsize\strut #1}}

\newcommand{\yes}{\textbf{x}}

\begin{table}[htbp]
\centering
\setlength{\tabcolsep}{2.5pt}
\renewcommand{\arraystretch}{1.25}

\resizebox{\textwidth}{!}{%
\begin{tabular}{|A|S||*{14}{C|}}
\hline
\multicolumn{2}{|c||}{\textbf{Studie}} &
\multicolumn{14}{c|}{\textbf{Konzepte}} \\
\hline

\multirow{2}{*}{\lhdr{4.3cm}{Autor(en)\\\& Jahr}} &
\multirow{2}{*}{\lhdr{2.8cm}{System}} &
\multicolumn{3}{c|}{\textbf{Ziel}} &
\multicolumn{3}{c|}{\textbf{Datenquellen}} &
\multicolumn{3}{c|}{\textbf{\ac{KI}-/\ac{LLM}-Ansatz}} &
\multicolumn{3}{c|}{\textbf{Evaluation}} &
\multicolumn{2}{c|}{\textbf{Kontext}} \\
\cline{3-16}

& &
\hdr{\ac{WCAG}-Konf.\ prüfen} &
\hdr{Barrieren korrigieren} &
\hdr{Entwickler unterstützen} &
\hdr{Tool-Reports (axe etc.)} &
\hdr{Interaktionsdaten (Playwright)} &
\hdr{Quellcode-Analyse} &
\hdr{\ac{LLM}-basiert} &
\hdr{Kontextsensitiv / \ac{RAG}} &
\hdr{Lokal ausführbar} &
\hdr{Automatisierte Eval.} &
\hdr{Qualitative Eval.} &
\hdr{Nutzerstudie} &
\hdr{Rich-Client / \ac{SPA}} &
\hdr{Rechtl.\ Rahmen (\ac{BFSG} etc.)} \\
\hline

% ── HOCH-priorisierte Studien ────────────────────────────────
\multicolumn{16}{|l|}{\textbf{Hoch priorisierte Studien}} \\
\hline

He et al.\ (2025) & GenA11y
& \yes &   &   & \yes &   &   & \yes &   &   & \yes &   &   &   &   \\ \hline

Fathallah et al.\ (2025) & AccessGuru
& \yes & \yes & \yes & \yes & \yes &   & \yes &   &   & \yes &   &   &   &   \\ \hline

Bassi et al.\ (2025) & \ac{LLM} Auditing
& \yes & \yes & \yes &   &   &   & \yes &   &   &   & \yes &   &   & \yes \\ \hline

Paternò et al.\ (2025) & TeVal
& \yes &   & \yes & \yes &   &   & \yes & \yes &   &   & \yes &   &   &   \\ \hline

Tafreshipour et al.\ (2024) & Ma11y
& \yes &   &   & \yes & \yes &   &   &   &   & \yes &   &   &   &   \\ \hline

Fernández-Navarro \& Chicano (2026) & \ac{LLM} Remediation
& \yes & \yes & \yes & \yes &   & \yes & \yes &   &   & \yes &   &   & \yes &   \\ \hline

% ── Eigener PoC ──────────────────────────────────────────────

\textbf{Martinez Campo (2026)} &
\textbf{ACE (\ac{PoC})}
& \yes & \yes & \yes
& \yes & \yes & \yes
& \yes & \yes & \yes
&   & \yes &
& \yes & \yes \\ \hline

\end{tabular}%
}% Ende resizebox

\caption[Konzeptmatrix (Hauptmatrix, hoch priorisierte Studien)]{%
  Konzeptmatrix nach Webster Watson\footnotemark: hoch priorisierte Studien im Vergleich
  zum eigenen \acf{PoC}. \yes~=~Konzept vorhanden/adressiert; leer ~=~nicht adressiert.
  Für den eigenen \ac{PoC} bezeichnen die Markierungen angestrebte Designziele, deren Umsetzung und Wirksamkeit in Kapitel~\ref{kap:7} evaluiert wird.
  Die vollständige Matrix aller analysierten Studien findet sich in Anhang~\ref{app:konzeptmatrix}.}
\label{tab:konzeptmatrix_haupt}
\end{table}
\footnotetext{Vgl. \cite{WebsterWatson2002}.}

\chapter{Methodik} \label{kap:methodik}

Das folgende Kapitel beschreibt das methodische Vorgehen der Arbeit. Zunächst wird die gewählte Forschungsmethode begründet (Abschnitt~\ref{sec:wiss_methode}), anschließend der Untersuchungsgegenstand abgegrenzt (Abschnitt~\ref{sec:scope}) und schließlich das Evaluationsdesign mit den zugehörigen Bewertungskriterien definiert (Abschnitt~\ref{sec:eval}).

\section{Wissenschaftliche Methode} \label{sec:wiss_methode}

Die vorliegende Arbeit folgt dem Forschungsparadigma des \ac{DSR}, das von Hevner et al. für die Wirtschaftsinformatik systematisiert wurde.\footcite[Vgl.][]{hevner2004dsr} \ac{DSR} eignet sich für diese Arbeit aus zwei Gründen: Erstens steht die Konstruktion eines konkreten technischen Artefakts im Zentrum der Forschungsfrage, nicht die Überprüfung einer statistischen Hypothese. Zweitens fordert \ac{DSR} explizit die Evaluation des Artefakts in einem realen Anwendungskontext, was dem praxisorientierten Charakter der Arbeit entspricht.\footcite[Vgl.][S. 80]{hevner2004dsr}

Alternative Ansätze wie Action Research oder kontrollierte Experimente wurden verworfen: Action Research setzt eine mehrjährige Praktikerzusammenarbeit voraus; kontrollierte Experimente erfordern eine Stichprobenbasis, die im Verwaltungskontext nicht verfügbar ist.\footcites[Vgl.][S. 1--2]{baskerville1999ar}[Vgl.][S. 239--335]{baskerville2004ar}

Als prozessualer Rahmen dient das sechsphasige DSRM-Prozessmodell nach Peffers et al.\footcite[Vgl.][]{peffers2007dsrm} Die Zuordnung der einzelnen Phasen zu den Kapiteln dieser Arbeit zeigt Tabelle~\ref{tab:dsrm_mapping}.

\begin{table}[htbp]
\centering
\begin{tabular}{|l|l|}
\hline
\textbf{DSRM-Phase} & \textbf{Kapitel} \\
\hline
Problem-Identifikation & Kapitel \ref{kap:1}-\ref{kap:3} \\
Zieldefinition & Kapitel \ref{kap:1.3} \& \ref{Forschungslücken} \\
Design \& Entwicklung & Kapitel \ref{kap:4}-\ref{kap:6} \\
Demonstration & Kapitel \ref{kap:6} (insb. Beispielabläufe) \\
Evaluation & Kapitel \ref{kap:7} \\
Kommunikation & Gesamtarbeit, zusammenfassend Kapitel \ref{kap:8} \\
\hline
\end{tabular}
\caption{Zuordnung der DSRM-Phasen zu den Kapiteln der Arbeit.}
\label{tab:dsrm_mapping}
\end{table}

Ergänzt wird dieser Rahmen durch einen Fallstudienansatz nach Yin.\footcite[Vgl.][]{yin_case_2018} Der \ac{BWEC} dient als Einzelfallstudie, da die Forschungsfrage auf ein zeitgenössisches Phänomen in einem spezifischen Praxiskontext abzielt, bei dem die Grenzen zwischen technischem System und organisatorischem Kontext nicht trennscharf sind.

Die Wahl eines \ac{PoC} als Artefakttyp ist durch die explorative Natur der Forschungsfrage begründet: Das Ziel ist die konzeptionelle Demonstration der technischen Machbarkeit, nicht die statistisch belastbare Generalisierung auf eine Grundgesamtheit.\footcite[Vgl.][]{hevner2004dsr}

\section{Untersuchungsgegenstand und Scope-Abgrenzung} \label{sec:scope}

Als Praxisrahmen und Motivationskontext dient der \ac{BWEC}, eine öffentlich zugängliche React-basierte Webanwendung der \ac{BITBW}.\footcite[Vgl.][]{bwec} Er vereint drei für die Forschungsfrage zentrale Merkmale: gesetzlich relevante Anforderungen an digitale Barrierefreiheit, einen datenschutzsensiblen Einsatzkontext in der öffentlichen Verwaltung sowie typische Interaktionsmuster moderner Single-Page Applications. Die formale Hauptevaluation erfolgt hingegen an synthetisch konstruierten Testsuiten (test-12, test-50, test-100) mit vollständig bekannter Ground Truth. Der \ac{BWEC} wird ergänzend als Feldtest ohne quantitativen Ground-Truth-Vergleich eingesetzt (vgl. Abschnitt~\ref{sec:eval_bwec}).

Der Untersuchungsscope ist wie folgt abgegrenzt:

\begin{itemize}
    \item \textbf{Eingeschlossen:} Öffentlich zugängliche Bereiche der Anwendung ohne Authentifizierung. Betrachtet werden \ac{WCAG}~2.2-Verstöße auf Konformitätsstufe A und AA, soweit sie durch regelbasierte Prüfung, \ac{DOM}-Analyse und Interaktionsausführung im Frontend beobachtbar sind.\footcite[Vgl.][]{WCAG}
    \item \textbf{Ausgeschlossen:} Backend-Logik, Datenbankstrukturen sowie eine vollständige rechtliche Bewertung im Sinne einer formalen \ac{BITV}-Prüfung.\footcite[Vgl.][]{BITV}
\end{itemize}

Die Einschränkungen hinsichtlich vollständiger \ac{WCAG}-Abdeckung und Generalisierbarkeit auf beliebige Frameworks werden in Kapitel~\ref{kap:7} als Limitationen diskutiert.

\section{Evaluationsdesign und Kriterien} \label{sec:eval}

Ziel der Evaluation ist es, die technische Machbarkeit und den praktischen Nutzen des entwickelten \ac{PoC} anhand synthetisch konstruierter Testsuiten mit vollständig bekannter Ground Truth zu untersuchen. Der \ac{BWEC} dient ergänzend als Feldtest-Kontext ohne quantitativen Recall-Vergleich (vgl. Abschnitt~\ref{sec:eval_bwec}).

\subsection{Evaluationslogik}

Die Evaluation folgt einem fünfstufigen Vorgehen:

\begin{enumerate}
    \item \textbf{Testsuite-Konstruktion:} Drei synthetische React-\acp{SPA} werden mit gezielt eingebetteten \ac{WCAG}-Violations entwickelt (test-12: 12, test-50: 50, test-100: 100~Violations), die alle drei Verstoßkategorien der in Abschnitt~\ref{Taxonomie} eingeführten Taxonomie abdecken.\footcite[Vgl.][]{WCAG}
    \item \textbf{Benchmarkläufe:} Jede Suite wird mit drei Modellen je zehn Mal analysiert (insgesamt 90~Runs). Pro Lauf werden Recall, Priorisierungsverhalten, Fix-Qualität und Laufzeit erfasst.
    \item \textbf{Recall-Auswertung:} Die Ergebnisse werden mit der bekannten Ground Truth verglichen. Eine Violation gilt als erkannt, wenn sie in der Mehrzahl der zehn Runs im finalen \ac{LLM}-Report erscheint.
    \item \textbf{Qualitative Empfehlungsbewertung:} Die generierten Handlungsempfehlungen werden anhand einer dreistufigen Bewertungsskala (Stufe~0--2) nach Dateikonkretheit und Code-Fix-Qualität bewertet.
    \item \textbf{\ac{BWEC}-Feldtest:} Die Pipeline wird ergänzend auf dem realen \ac{BWEC}-Untersuchungsobjekt eingesetzt, um Systemstabilität und qualitative Plausibilität unter realen Bedingungen zu prüfen.\footcites[Vgl.][]{axe}[Vgl.][]{kodithuwakku_assessing_2025}
\end{enumerate}

\subsection{Bewertungskriterien}

Die Evaluation erfolgt anhand von vier Kriterien, die auf Basis der Forschungsfragen und der Fachliteratur entwickelt wurden.

\textbf{Kriterium 1 -- Korrektheit der Erkennung:} Die vom System identifizierten Verstöße werden mit der bekannten Ground Truth der synthetischen Testsuiten verglichen. Als Orientierungswerte dienen die Ergebnisse vergleichbarer Systeme: \textfett{GenA11y} erreicht eine Precision von 94,5\% und einen Recall von 87,61\%.\footcite[Vgl.][FSE101:20]{he_enhancing_2025}

\textbf{Kriterium 2 -- Qualität der Handlungsempfehlungen:} Die generierten Empfehlungen werden anhand einer dreistufigen Skala (Stufe~0--2) qualitativ bewertet: Stufe~2 bezeichnet einen konkreten Fix mit Dateiname, Zeilennummer und umsetzbarem Code-Vorschlag; Stufe~1 einen zutreffenden Fix-Typ ohne Quellcode-Kontext; Stufe~0 eine generische Problembeschreibung. \textfett{AccessGuru} erreicht als Referenzwert einen Violation Score Decrease von 84\%.\footcite[Vgl.][]{fathallah_accessguru_2025}

\textbf{Kriterium 3 -- Mehrwert der Kontextualisierung:} Der Zusatznutzen der vollständigen Kontextkombination wird durch einen systematischen Vergleich zweier Konfigurationen ermittelt: (a)~die reine Tool-Baseline (axe-core, Playwright, grep ohne \ac{LLM}-Kontextualisierung) und (b)~die vollständige Pipeline (Baseline + \ac{LLM}-gestützte Quellcodeanalyse). Verglichen werden Recall, Fix-Qualität sowie die Fähigkeit, datei- und zeilenspezifische Handlungsempfehlungen zu generieren.\footnote{Siehe Abschnitt~\ref{kap:1.3}.}

\textbf{Kriterium 4 -- Technische Praktikabilität:} Laufzeit, Stabilität und Qualität des \acs{JSON}-Outputs werden unter realen Bedingungen dokumentiert.

\subsection{Limitationen des Evaluationsdesigns}

Zwei methodische Einschränkungen sind vorab zu benennen. Erstens wird die Bewertung durch den Autor selbst durchgeführt und nicht durch unabhängige Gutachter validiert. Zweitens sind die genannten Referenzwerte von \textfett{AccessGuru}\footcite[Vgl.][]{fathallah_accessguru_2025} und \textfett{GenA11y}\footcite[Vgl.][]{he_enhancing_2025} nur eingeschränkt vergleichbar, da sie auf anderen Datensätzen und unter anderen Rahmenbedingungen ermittelt wurden. Sie dienen daher als Orientierung, nicht als direkter Benchmark.